\section{HTTP Interface}
\label{sec:http}

The browser interface talks to \texttt{fsb} over HTTP.  This interface is described
so that it can be tested and reviewed.  It is \emph{internal}: it may change between
versions, it is not a supported programming interface, and it is refused to anything
that does not hold the session (Section~\ref{sec:security}).

All paths in this section are relative to the random prefix (\rid{SEC-7}): the
interface is at \code{/PREFIX/}, and the endpoints at \code{/PREFIX/api/\ldots}.  Every
endpoint takes the request path in the query parameter \code{path}, an absolute path,
decoded once; a literal \code{\%2e\%2e} is just part of a name.

\subsection{Endpoints}

\begin{req}{API-1}
The interface answers the endpoints of the following table, and no others.  All are
\code{GET} (or \code{HEAD}).

\begin{center}
\small
\begin{tabular}{@{}p{0.17\linewidth}p{0.20\linewidth}p{0.55\linewidth}@{}}
\toprule
Endpoint & Parameters & Answer \\
\midrule
\code{/} & & The interface (static files). \\
\code{api/status} & & JSON: \code{readOnly} (always true), \code{coreDenyMissing} (list), \code{roots}, \code{home}. \\
\code{api/list} & \code{path} & The entries of a folder, streamed (\rid{API-2}). \\
\code{api/file} & \code{path} & The file's bytes as a download (\rid{API-4}). \\
\code{api/head} & \code{path}, \code{bytes} & A classified look at the start of a file (\rid{API-5}). \\
\code{api/meta} & \code{path}, \code{values} & Details and extended attributes (\rid{API-6}). \\
\code{api/preview} & \code{path} & A raster image for inline display. \\
\code{api/pdf} & \code{path} & A PDF for inline display. \\
\code{api/archive} & \code{path} & The table of contents of an archive. \\
\code{api/quicklook} & \code{path} & The Quick Look picture of an iWork or Office document (\rid{API-9}). \\
\code{api/search} & \code{path}, \code{q}, \code{case} & Name matches under a folder, streamed (\rid{API-3}). \\
\code{api/mdframe} & & The fixed page that renders Markdown. \\
\bottomrule
\end{tabular}
\end{center}
\end{req}

\begin{req}{API-2}
\code{api/list} answers with newline-delimited JSON: a line \code{\{"path":P\}}, then
lines \code{\{"entries":[\ldots]\}} in the order the file system gave them (the client
sorts), then \code{\{"done":true\}}, or, if the listing broke after it had started,
\code{\{"error":"listing interrupted"\}} instead.  A client treats a stream that ends
with neither as broken, never as a complete listing.  An entry has \code{name}, \code{isDir},
\code{size}, \code{modTime} and \code{mode}, and, when true, \code{isSymlink} and
\code{broken}.  A denied, hidden or unreachable entry is not present.  A path that is
not a folder, is denied, or does not exist is answered 404.
\end{req}

\begin{req}{API-3}
\code{api/search} answers with newline-delimited JSON: a line
\code{\{"path":P,"query":Q\}}, then lines \code{\{"matches":[\ldots]\}} as results are
found (each with \code{path}, \code{rel}, \code{name}, \code{isDir}, \code{size} and
\code{modTime}), then \code{\{"done":true,"visited":N,"truncated":B,"incomplete":B\}}, or
\code{\{"error":"search interrupted"\}} instead.  \code{incomplete} is true when part
of what the search was allowed to look at could not be searched (\rid{SRC-6}); the
interface then says so.  A client treats a stream that ends with neither line as
broken.  \code{case=1} makes matching case-sensitive.
The limits are those of \rid{SRC-2}.  An empty query finds nothing.
\end{req}

\begin{req}{API-4}
\code{api/file} answers a regular file with \code{Content-Type:
application/octet-stream} and \code{Content-Disposition: attachment; filename=\ldots}
and supports range requests (status 206) for large files.  A folder is answered 404.
\end{req}

\begin{req}{API-5}
\code{api/head} reads at most \code{bytes} bytes (default 2048, at most 65{,}536; a
value outside 1 to 65{,}536 is answered 400) and answers JSON with \code{kind}
(\code{text}, \code{binary}, \code{empty} or \code{dataless}), \code{size}, and, for
text, \code{text} and \code{truncated}.  Text is returned only if it is valid UTF-8; a
binary or cloud-only file's bytes are never returned.
\end{req}

\begin{req}{API-6}
\code{api/meta} answers JSON with \code{name}, \code{path}, \code{isDir}, \code{size},
\code{modTime}, \code{mode}, \code{isSymlink}, \code{symlinkTarget}, \code{dataless}
and \code{xattrs} (each with \code{name}, and, unless \code{values=0}, \code{value},
\code{encoding} (\code{utf8} or \code{hex}), \code{size} and \code{large}).
\end{req}

\begin{req}{API-7}
\code{api/preview}, \code{api/pdf} and \code{api/archive} answer only for a file whose
bytes are of the corresponding kind (\rid{PRV-2}); any other file is answered 415.
\code{api/preview} and \code{api/pdf} support range requests.
\end{req}

\begin{req}{API-9}
\code{api/quicklook} answers a PNG picture, sent like \code{api/preview} and never
cached, for a document of one of the types of \rid{PRV-12}, file or package.  It is
answered 415 for any other type, for a document larger than 256~MB, and when Quick
Look cannot draw the document in time; 409 for a cloud-only document or package part;
and 404 for a denied or missing document.  A package answers the same whether or not
it holds parts that a listing leaves out.  It does not support range requests.  Quick
Look is part of macOS: on any other system every request is answered 415 at once,
without copying the document or starting any program.
\end{req}

\begin{req}{API-10}
Every name and path in a JSON answer (the entries of \code{api/list} and
\code{api/search}, the \code{name}, \code{path}, \code{symlinkTarget} and attribute
names of \code{api/meta}, and the \code{roots} and \code{home} of \code{api/status})
keeps all of its bytes.  (The members of an archive are names inside a file, only ever
displayed; they are cleaned as \rid{PRV-9} describes.)  A name that is valid UTF-8 is sent as it
is.  In one that is not (Linux allows any bytes in a name but \code{/} and NUL), each
byte that is not part of valid UTF-8 is sent as a NUL character followed by its value
in two uppercase hexadecimal digits, so \code{caf}\textit{E9}\code{.txt} is sent as
\code{"caf\textbackslash u0000E9.txt"}.  Since no name contains NUL, the escape is
never ambiguous.  The client sends such a path back with each escape replaced by the
byte itself (\code{caf\%E9.txt}), so the \code{path} parameter always carries the
name's real bytes; the server does not decode the escape.
\end{req}

\subsection{Status codes}

\begin{req}{API-8}
The status codes and their meanings are the following.

\begin{center}
\small
\begin{tabular}{@{}lp{0.75\linewidth}@{}}
\toprule
Status & Meaning \\
\midrule
200, 206 & Success (206: a range of the file). \\
303 & The launch token was accepted; the session cookie is set (\rid{SEC-5}). \\
400 & A parameter is out of range (\code{bytes}). \\
403 & No valid session, a wrong \code{Host}, a cross-origin request, a URL outside the prefix, or, on Linux, a launch or a request with the session from a connection not of \texttt{fsb}'s own user, or whose owner cannot be found (Section~\ref{sec:security}, \rid{SEC-16}); or a permission error on a non-denied path (\rid{RUL-7}). \\
404 & Not found: the path does not exist, is denied, lies outside the roots, or is not a regular file or folder (\rid{SEC-14}). \\
405 & A method other than \code{GET} or \code{HEAD}. \\
409 & The file is stored only in the cloud and was not read (\rid{SEC-13}). \\
415 & The file is not of the kind the endpoint serves (\rid{API-7}). \\
416 & A range that cannot be satisfied. \\
500 & An internal error (the details are logged, not sent). \\
\bottomrule
\end{tabular}
\end{center}
\end{req}
